\Needspace{28\baselineskip}
\begin{center}
PALLAS.
\end{center}


Monatliche Correspondenz. Band VII, S.\ 556—560. Juni 1803.


— — —

Dr.\ \textsc{Gauß} vergleicht unermüdlich weiter die Beobachtungen der \emph{Pallas} von \textsc{Olbers} mit seinen Elementen VI. Hier ist das Ergebnis der Beobachtungen vom 16.\ März bis zum 13.\ April. Man wird mit Vergnügen sehen, dass diese Elemente weiterhin gut übereinstimmen und die Unterschiede kaum größer sind als die bei den Beobachtungen selbst zu erwartenden Unsicherheiten, besonders bei den Deklinationen, die Dr.\ \textsc{Olbers} selbst als sehr unsicher angibt:

\begin{center}
\begingroup\fontsize{8}{10}\selectfont\setlength{\tabcolsep}{3pt}\renewcommand{\arraystretch}{1.12}
\begin{tabular}{r|r|r|r|r|r}
1803 & \shortstack{Mittlere Zeit\\in\\Bremen} & \shortstack{Berechnete\\Rektaszension\\der Pallas} & \shortstack{Berechnete\\Norddeklination\\der Pallas} & \shortstack{Unterschied\\in Rektasz.\\R—B} & \shortstack{in Dekl.\\R—B} \\ \hline
März 16 & $14^u\,10^m\,28^s$ & $278^\circ\,37'\,10''1$ & $10^\circ\,43'\,18''7$ & $-5''9$ & $+57''7$ \\
21 & $13\,40\,47$ & $279\,37\,19.0$ & $11\,30\,22.1$ & $+14.0$ & $+19.1$ \\
22 & $13\,14\,22$ & $279\,48\,30.2$ & $11\,39\,49.6$ & $+20.2$ & $+64.6$ \\
24 & $13\,2\,12$ & $280\,10\,29.9$ & $11\,59\,11.0$ & $+4.9$ & $+28.0$ \\
31 & $13\,24\,25$ & $281\,20\,15.1$ & $13\,8\,54.8$ & $-2.9$ & $+4.8$ \\
April 11 & $12\,10\,9$ & $282\,43\,33.5$ & $15\,0\,59.4$ & $-16.5$ & $+9.4$ \\
12 & $12\,12\,18$ & $282\,49\,32.1$ & $15\,11\,21.3$ & $-11.9$ & $+4.3$ \\
13 & $12\,13\,41$ & $282\,55\,13.0$ & $15\,21\,42.8$ & $-2.0$ & $+24.8$ \\
\end{tabular}\endgroup
\end{center}

\begin{flushright}
[\textsc{von Zach}.]
\end{flushright}
